% ************************ Informações do Template e Meta-Data **********************
%% Título do Trabalho

\newcommand{\HRule}{\rule{\linewidth}{0.5mm}}



\newcommand{\TITULOa}{\textbf{Estimativas de curvatura para sólitons de Ricci gradientes quadri-dimensionais e classificação de quase-sólitons de Ricci gradientes com curvatura harmônica}}



\newcommand{\TITULOb}{\textbf{Sobre sólitons de Ricci gradientes shrinking quadri-dimensionais e quase-sólitons de Ricci gradientes com curvatura harmônica}}


\newcommand{\TITULOc}{\textbf{Sobre estimativas de curvatura para sólitons de Ricci gradientes shrinking quadri-dimensionais e estrutura local de quase-sólitons de Ricci gradientes com curvatura harmônica}}

\newcommand{\TITULOfinal}{\textbf{Título teste a ser usado para teste e preenchimento do template}}


\title{\TITULOfinal}






%\HRule \\[0.2cm]
%{ \huge \textit{Uma (breve) introdução às variedades diferenciáveis} \\[0.2cm] % Title of your document
%\HRule \\[1.5cm]

 %\texorpdfstring{\\ \LaTeX2e}{LaTeX2e}}
%\texorpdfstring é usado para informações de metadata do PDF - use para simplificar símbolos do título na metadata
%\texorpdfstring{LaTeX_Version}{PDF Version (non-latex)} eg.,
%\texorpdfstring{$sigma$}{sigma}

%% Subtítulo (opcional) - comente esta linha se o trabalho não possuir subtítulo
%\subtitle{Subtítulo}

%% Nome Completo do Autor
\author{Fulano da Silva}



%% Departamento ()
%\dept{Department of Mathematics}
\dept{Departamento de Matemática,
}

%% Universidade
\university{Universidade de Brasília}
\crest{\includegraphics[width=1\textwidth]{unb.eps}}


%% Não Altere esta linha
\collegeshield{\includegraphics[width=0.2\textwidth]{CollegeShields/empty_shield.png}}


%% Orientador (opcional)
%% para mútiplos orientadores separe-os com o comando \newline
%\supervisor{Prof. A.B. Co-orientador\newline
%Prof. C.D. Orientador}

%% Nome (na capa) do co-orientador - Supervisor (default) ou deixe abaixo descomentado se seu trabalho for escrito em Português
%% Se o trabalho for em Inglês, simplemente comente a linha abaixo
%\supervisorrole{Co-Orientador: }
%% if no title is desired:
% \supervisorrole{}

%% Supervisor line width: required to align supervisors
%\supervisorlinewidth{0.50\textwidth}

%% Advisor (optional)
%% for multiple advisors, append each advisor with the \newline command
\advisor{Beltrano de Souza}

%% Advisor Role (optional) - Advisor (default) or leave empty
\advisorrole{Orientador: }
%% if no title is required
% \advisorrole{}

%% Comprimento da linha do nome do Orientador - altere se ficar desalinhado (seu orientador tem um nome muito grande ?)
%\advisorlinewidth{0.50\textwidth}


%% Texto de Submissão, deixe descomentado de acordo com a língua que está escrevendo.
%% Português
\renewcommand{\submissiontext}{Dissertação apresentada como requisito parcial para obtenção do grau de}

%% Inglês
%\renewcommand{\submissiontext}{Dissertation submitted in partial fulfillment of the requirements for the degree of}

%% Nome completo do Grau - Mestre em Matemática (Doutor em Matemática)
\degreetitle{Mestre em Matemática}

%% Afiliação do Departamento (opcional)
%\college{Instituto de Ciências Exatas}

%% Submission date
% Default is set as {\monthname[\the\month]\space\the\year}
\degreedate{Brasília, DD de MMMMM de 2025}

%% Meta informação - palavras chave de seu trabalho
\subject{Matemática} \keywords{{fluxo de Ricci}, {sólitons de Ricci}, {estimativas de curvatura}, {tensor de Einstein}, {variedades quadridimensionais.}}
